% 第十三节：焦点与准线
\chapter{焦点与准线}

\section{复点与垂直结构}

平行与正交一章说明，正交结构可以看作由两个无穷远处的复点确定。记
\begin{equation}
z^i=(1,\ii,0),\qquad \bar z^i=(1,-\ii,0),
\label{eq:zz}
\end{equation}
其中 $\ii$ 为虚数单位。两点都在无穷远直线 $\linf$ 上，事实上
\[
z\times\bar z=(0,0,-2\ii)\sim(0,0,1)=\linf .
\]
$z$ 与 $\bar z$ 互为复共轭，它们与实射影平面上的实点、实直线具有相同的代数地位。

本节的约定是：把坐标含复数的点矢、线矢与实的点矢、线矢同样处理，一切运算（线性组合、缩并、
叉乘、求逆）都按代数规则进行，只在需要判断某个量是否为实数时才使用复共轭。经验表明，只要问题
的提法本身不涉及实数性，坐标中的虚部不影响结论；下文凡按几何意义应当为实的量，算出的结果都是
实的。

垂直结构取为这两点的对称积
\begin{equation}
\theta^{ij}=(\sym{z}{\bar z})^{ij}=z^i\bar z^j+z^j\bar z^i .
\label{eq:theta-zz}
\end{equation}
按分量计算，$z^i\bar z^j$ 只在 $(1,1)$ 与 $(2,2)$ 处分别为 $1$ 与 $1$，其余分量为零，因此
\[
\theta^{ij}=
\begin{pmatrix}
2&0&0\\
0&2&0\\
0&0&0
\end{pmatrix}
\sim\diag\{1,1,0\}.
\]
可见 \eqref{eq:theta-zz} 定义的 $\theta$ 与实标准形式 $\diag\{1,1,0\}$ 只相差规范因子 $2$，属于
同一个倍数类。在下文的分量计算中直接取
\begin{equation}
\theta^{ij}=\diag\{1,1,0\}.
\label{eq:theta-standard}
\end{equation}

对任意线矢 $\ell$，
\begin{equation}
\theta^{ij}\ell_i\ell_j=2(\ell\cdot z)(\ell\cdot\bar z),
\label{eq:theta-line}
\end{equation}
因此 $\theta$ 所对应的退化二次点正是两点 $z,\bar z$：经过 $\theta$ 的直线恰为过 $z$ 的直线与过
$\bar z$ 的直线的并。$\theta$ 的号差为 $(1,1,0)$，它是退化的，核由 $\linf$ 张成。

\section{焦点方程与离心率}

\begin{definition}[焦点方程]
设 $g$ 是一条圆锥曲线，$\theta$ 是垂直结构。方程
\begin{equation}
\det\bigl(\theta^{ij}-\lambda g^{ij}\bigr)=0
\label{eq:focus-eq}
\end{equation}
称为 $g$ 的\textbf{焦点方程}，其中的 $\lambda$ 称为\textbf{特征值}。这里 $g^{ij}$ 是 $g_{ij}$ 的
逆 $(g^{-1})^{ij}$。
\end{definition}

焦点方程是三维矢量空间中的特征值方程：把 $\theta,g$ 看作矩阵，由 $g$ 可逆，
\[
\det(\theta-\lambda g)=\det g\cdot\det\bigl(g^{-1}\theta-\lambda I\bigr),
\]
故 \eqref{eq:focus-eq} 的根正是线性变换 $g^{-1}\theta$ 的特征值。由于 $g^{-1}\theta$ 是三阶
矩阵，方程关于 $\lambda$ 至多是三次的。又 $\theta$ 退化，$\det\theta=0$，即
$\det(g^{-1}\theta)=0$，所以 $\lambda=0$ 是一个根。

\begin{proposition}[特征值]
焦点方程 \eqref{eq:focus-eq} 有三个根
\[
\lambda_1,\qquad \lambda_2,\qquad 0,
\]
其中 $\lambda_1,\lambda_2$ 是二次方程
\[
\det\bigl(\theta^{ij}-\lambda g^{ij}\bigr)=0
\]
在去掉因子 $\lambda$ 之后所剩的两个根。
\end{proposition}

\begin{proposition}[规范变换下特征值的变化]
把 $\theta$ 换成 $\nu\theta$、把 $g$ 换成 $\mu g$（$\mu,\nu\ne0$）时，\eqref{eq:focus-eq} 的
非零根全部乘以同一个因子 $\nu/\mu$。
\end{proposition}

\begin{proof}
直接计算
\[
\det(\nu\theta-\lambda\mu g)
=\nu^3\det\Bigl(\theta-\frac{\lambda\mu}{\nu}g\Bigr).
\]
因此 $\nu\theta-\lambda\mu g$ 退化当且仅当 $\theta-\lambda_0g$ 退化，其中
$\lambda_0=\lambda\mu/\nu$，即 $\lambda=\lambda_0\nu/\mu$。原来的每个特征值都乘以 $\nu/\mu$。
\end{proof}

由此，单个特征值 $\lambda_1,\lambda_2$ 依赖 $g,\theta$ 的规范选取，只有比值
$\lambda_2/\lambda_1$ 是规范无关的，因而有几何意义。据此定义离心率。

\begin{definition}[离心率]
在焦点方程的两个非零根中，约定 $\lambda_1$ 对应实焦点所在的那条主轴（见下一节），$\lambda_2$
对应另一条主轴，并取
\begin{equation}
e=\sqrt{1-\frac{\lambda_2}{\lambda_1}}.
\label{eq:ecc}
\end{equation}
称 $e$ 为圆锥曲线 $g$ 的\textbf{离心率}。
\end{definition}

\begin{remark}[次序的约定]
$\lambda_1,\lambda_2$ 是两个根，本身没有先后之分；\eqref{eq:ecc} 只在固定取法之后才有确定的值。
上文把 $\lambda_1$ 约定为对应实焦点所在主轴的根，这使 $e$ 在下述意义下取到标准值：椭圆时
$0<e<1$，抛物线时 $e=1$，双曲线时 $e>1$，圆时 $e=0$。若不约定次序，则 $e$ 可能对应
$1-e^2$ 的另一个分支。
\end{remark}

\section{六个焦点与三条主轴}

由上一节的结论，$\lambda_1,\lambda_2,0$ 分别对应三个退化二次点
\begin{equation}
X^{ij}=\theta^{ij}-\lambda_1g^{ij},\qquad
Y^{ij}=\theta^{ij}-\lambda_2g^{ij},\qquad
\theta^{ij}=\theta^{ij}-0\cdot g^{ij}.
\label{eq:XYZ}
\end{equation}
三者都满足 $\det=0$。一般情况下它们的秩为 $2$，号差为 $(1,1,0)$，因而每一个都可以写成两个点的
对称积。按第八章的记号，设
\begin{equation}
X=(\sym{f_{X1}}{f_{X2}}),\qquad
Y=(\sym{f_{Y1}}{f_{Y2}}),\qquad
\theta=(\sym{z}{\bar z}).
\label{eq:focus-sym}
\end{equation}

\begin{definition}[焦点]
\eqref{eq:focus-sym} 中的六个点
\[
f_{X1},\ f_{X2},\ f_{Y1},\ f_{Y2},\ z,\ \bar z
\]
称为圆锥曲线 $g$ 的\textbf{焦点}。其中 $f_{X1},f_{X2}$ 称为一对配对焦点，$f_{Y1},f_{Y2}$ 与
$z,\bar z$ 也各称一对配对焦点；每一对配对焦点的对称积依次为 $X,Y,\theta$。
\end{definition}

于是圆锥曲线共有六个焦点，两两配对成三对。

\begin{proposition}[主轴]
设 $x_i$ 是 $\lambda_1$ 对应的特征线矢，即
$X^{ij}x_j=0$，则直线 $x$ 经过 $f_{X1},f_{X2}$；同理 $\lambda_2$ 对应的特征线矢 $y$ 经过
$f_{Y1},f_{Y2}$，而 $\lambda=0$ 对应的特征线矢是无穷远直线 $\linf$。这三条直线分别称为
$x$ \textbf{主轴}、$y$ \textbf{主轴}与无穷远直线。
\end{proposition}

\begin{proof}
由 $X^{ij}=f_{X1}^if_{X2}^j+f_{X2}^if_{X1}^j$，对任意线矢 $x$，
\[
X^{ij}x_j=(f_{X2}\cdot x)f_{X1}^i+(f_{X1}\cdot x)f_{X2}^i,
\]
其中 $f\cdot x=f^ix_i$。$f_{X1},f_{X2}$ 不共线，故 $X^{ij}x_j=0$ 当且仅当
$f_{X1}\cdot x=f_{X2}\cdot x=0$，即直线 $x$ 同时经过 $f_{X1},f_{X2}$。对 $\lambda_2$ 同理。

对 $\lambda=0$，条件 $\theta^{ij}\linfsub{j}=0$ 由 \eqref{eq:theta-standard} 直接看出，核由
$\linf$ 张成。故 $\lambda=0$ 对应的特征线矢就是 $\linf$。
\end{proof}

\begin{proposition}[主轴两两共轭]
$x,y,\linf$ 三条直线两两关于 $g$ 共轭，即对其中任意两条 $m,n$ 都有
$g^{ij}m_in_j=0$。
\end{proposition}

\begin{proof}
设 $m,n$ 分别对应特征值 $\lambda,\mu$，即
$\theta^{ij}m_j=\lambda g^{ij}m_j$、$\theta^{ij}n_j=\mu g^{ij}n_j$。用 $n_i$ 缩并第一式，
用 $m_i$ 缩并第二式，由于 $\theta^{ij}$ 对称，两式左端相等，于是
\[
\lambda\,g^{ij}m_in_j=\mu\,g^{ij}m_in_j,
\qquad\text{即}\qquad
(\lambda-\mu)\,g^{ij}m_in_j=0 .
\]
当 $\lambda\ne\mu$ 时得 $g^{ij}m_in_j=0$。三个特征值 $\lambda_1,\lambda_2,0$ 两两不同时结论
成立；若有重根，则该情形对应圆或抛物线的退化位置，此时相应的两条主轴重合或其中之一退化，
命题按极限理解。
\end{proof}

\section{圆锥曲线的分类}

特征值的符号决定圆锥曲线的类型。

\begin{proposition}[按特征值分类]
设 $\lambda_1,\lambda_2$ 是两个非零特征值，按上一节的约定排列。
\begin{enumerate}
\item 若 $\lambda_1\lambda_2>0$，则 $g$ 是\textbf{椭圆}，此时 $\lambda_2/\lambda_1\in(0,1)$，
$0<e<1$；
\item 若 $\lambda_1\lambda_2=0$（即某个非零根为零），则 $g$ 是\textbf{抛物线}，$e=1$；
\item 若 $\lambda_1\lambda_2<0$，则 $g$ 是\textbf{双曲线}，$e>1$。
\end{enumerate}
\end{proposition}

两类特殊的圆锥曲线也由特征值刻画。

\begin{proposition}[圆与等轴双曲线]
\begin{enumerate}
\item $g$ 是圆，当且仅当焦点方程有重根，即 $\lambda_1=\lambda_2$，此时 $e=0$；
\item $g$ 是等轴双曲线，当且仅当 $\lambda_1+\lambda_2=0$，此时 $e=\sqrt2$。
\end{enumerate}
\end{proposition}

\begin{proof}
第二款：由第一款的注记，若把 $\theta$ 降指标，则
\[
\lambda_1+\lambda_2=\tr(g^{-1}\theta)=g^{ij}\theta_{ij}=g_{ij}\theta^{ij}=g:\theta .
\]
等轴双曲线的定义是两条渐近线互相垂直，在正交几何中等价于 $g:\theta=0$，故
$g$ 为等轴双曲线当且仅当 $\lambda_1+\lambda_2=0$。代入 \eqref{eq:ecc} 得
$e=\sqrt{1-(-1)}=\sqrt2$。
\end{proof}

\begin{remark}[特征值的规范依赖]
$\lambda_1,\lambda_2$ 各自随规范改变，但 $\lambda_1=\lambda_2$、$\lambda_1\lambda_2=0$ 与
$\lambda_1+\lambda_2=0$ 都是规范无关的条件，因为规范变化把两者乘以同一个因子。因此上面的分类
条件有意义。
\end{remark}

\section{椭圆情形的分量计算}

下面取椭圆具体算出全部焦点，以验证上面的定义与解析几何中的焦点一致。设椭圆为
\[
\frac{x^2}{a^2}+\frac{y^2}{b^2}=1,\qquad a>b>0,
\]
其张量及其逆为
\begin{equation}
g_{ij}=\diag\Bigl\{\frac{1}{a^2},\frac{1}{b^2},-1\Bigr\},
\qquad
g^{ij}=\diag\{a^2,b^2,-1\}.
\label{eq:ellipse-g}
\end{equation}
按 \eqref{eq:theta-standard} 取 $\theta^{ij}=\diag\{1,1,0\}$，则
\[
\theta^{ij}-\lambda g^{ij}
=\diag\{1-\lambda a^2,\ 1-\lambda b^2,\ \lambda\},
\]
于是
\begin{equation}
\det\bigl(\theta^{ij}-\lambda g^{ij}\bigr)
=\lambda\,(1-\lambda a^2)(1-\lambda b^2).
\label{eq:ellipse-det}
\end{equation}
三个根为
\[
\lambda=0,\qquad \lambda=\frac{1}{a^2},\qquad \lambda=\frac{1}{b^2}.
\]
因 $a>b>0$，有 $1/b^2>1/a^2$。按约定 $\lambda_1$ 对应实焦点所在主轴，取
\begin{equation}
\lambda_1=\frac{1}{b^2},\qquad \lambda_2=\frac{1}{a^2},
\qquad
\frac{\lambda_2}{\lambda_1}=\frac{b^2}{a^2}.
\label{eq:ellipse-lambda}
\end{equation}
记 $c^2=a^2-b^2>0$，则
\[
e=\sqrt{1-\frac{\lambda_2}{\lambda_1}}
=\sqrt{1-\frac{b^2}{a^2}}
=\frac{c}{a},
\]
与解析几何中的离心率一致。

把 $\lambda_1=1/b^2$ 代入 \eqref{eq:XYZ}，得
\[
X^{ij}=\theta^{ij}-\frac{1}{b^2}g^{ij}
=\diag\Bigl\{-\frac{c^2}{b^2},\ 0,\ \frac{1}{b^2}\Bigr\}
\sim\diag\{-2c^2,\ 0,\ 2\}.
\]
后一个矩阵可以分解为对称积
\begin{equation}
X^{ij}=
\begin{pmatrix}
-2c^2&0&0\\
0&0&0\\
0&0&2
\end{pmatrix}
=
\begin{pmatrix}
-c\\0\\1
\end{pmatrix}
\begin{pmatrix}
-c&0&1
\end{pmatrix}
+
\begin{pmatrix}
c\\0\\1
\end{pmatrix}
\begin{pmatrix}
c&0&1
\end{pmatrix}.
\label{eq:ellipse-X}
\end{equation}
即
\[
X=(\sym{f_{X1}}{f_{X2}}),\qquad
f_{X1}=(-c,0,1),\qquad f_{X2}=(c,0,1).
\]
在仿射图 $z=1$ 中，这一对焦点的坐标是 $(-c,0)$ 与 $(c,0)$，正是解析几何中椭圆的两个焦点。

把 $\lambda_2=1/a^2$ 代入 \eqref{eq:XYZ}，得
\[
Y^{ij}=\theta^{ij}-\frac{1}{a^2}g^{ij}
=\diag\Bigl\{0,\ \frac{c^2}{a^2},\ \frac{1}{a^2}\Bigr\}
\sim\diag\{0,\ 2c^2,\ 2\},
\]
它同样可以分解为对称积
\begin{equation}
Y^{ij}=
\begin{pmatrix}
0&0&0\\
0&2c^2&0\\
0&0&2
\end{pmatrix}
=
\begin{pmatrix}
0\\ \ii c\\1
\end{pmatrix}
\begin{pmatrix}
0&-\ii c&1
\end{pmatrix}
+
\begin{pmatrix}
0\\ -\ii c\\1
\end{pmatrix}
\begin{pmatrix}
0&\ii c&1
\end{pmatrix}.
\label{eq:ellipse-Y}
\end{equation}
因此
\[
f_{Y1}=(0,\ii c,1),\qquad f_{Y2}=(0,-\ii c,1),
\]
它们的仿射坐标是 $(0,\ii c)$ 与 $(0,-\ii c)$，不在实平面内。

最后，$\lambda=0$ 对应的一对焦点是 $z=(1,\ii,0)$ 与 $\bar z=(1,-\ii,0)$，它们位于无穷远
直线上。

\begin{remark}[解析几何中可见的焦点]
六个焦点中，$f_{Y1},f_{Y2}$ 的坐标含虚数，$z,\bar z$ 在无穷远，只有 $f_{X1},f_{X2}$ 落在实
仿射平面内。因此在解析几何中通常只讨论 $(-c,0)$ 与 $(c,0)$ 这两个焦点；另外四个焦点是射影
构造中不可省去的部分，但在实平面内看不到。
\end{remark}

此时三条主轴分别是
\[
x=(0,1,0)\quad(y=0),\qquad
y=(1,0,0)\quad(x=0),\qquad
\linf=(0,0,1).
\]
$x$ 主轴是 $f_{X1},f_{X2}$ 的连线，即 $x$ 轴；$y$ 主轴是 $f_{Y1},f_{Y2}$ 的连线，即 $y$ 轴。
用 $g^{ij}=\diag\{a^2,b^2,-1\}$ 直接验证三者两两共轭：
$g^{ij}x_iy_j=g^{ij}x_i\linfsub{j}=g^{ij}y_i\linfsub{j}=0$，因为 $g^{ij}$ 是对角矩阵。

\section{六个焦点的几何构造}

现在说明六个焦点的几何来源。从 \eqref{eq:focus-eq} 出发，考察曲线系
\[
\theta^{ij}-\lambda g^{ij}.
\]
其中的成员都是 $(2,0)$ 型对称张量，即二次点。$\theta$ 与 $g^{ij}$ 各自有明确的几何意义：
\begin{itemize}
\item $\theta$ 是退化二次点 $(\sym{z}{\bar z})$。由 \eqref{eq:theta-line}，经过 $\theta$ 的直线恰
为过 $z$ 或过 $\bar z$ 的直线；
\item $g^{ij}$ 是圆锥曲线 $g$ 的对偶二次点，经过它的直线恰为 $g$ 的全部切线。
\end{itemize}
因此曲线系 $\theta-\lambda g$ 的每个成员都含有 $\theta$ 与 $g$ 共同含有的直线，也就是既经过
$z$ 或 $\bar z$、又与 $g$ 相切的直线。这样的直线共有四条：从 $z$ 引出的两条切线与从 $\bar z$
引出的两条切线。记这四条切线为 $\ell_1,\ell_2,\ell_3,\ell_4$。

另一方面，$X,Y,\theta$ 都是曲线系中行列式为零的成员，即退化二次点。一个秩为 $2$ 的退化二次点
对应两个点，经过它的直线恰为过其中任一者的直线。既然 $X$ 含有全部四条切线
$\ell_1,\ell_2,\ell_3,\ell_4$，这四条直线就必须分成两对，每一对各自过 $X$ 的两个点之一。把四条
直线分成两对的方式共有三种；三种分法分别给出 $X,Y,\theta$ 三个退化二次点。$\theta$ 对应的分法
是把过 $z$ 的两条切线放在一起、过 $\bar z$ 的两条放在一起，两个点正是 $z,\bar z$；$X$ 与 $Y$
对应的分法给出其余四个焦点。

把四条切线两两求交，共得 $6$ 个交点。上述三种分法所对应的恰好是这些交点的三种配对，六个焦点
正是这 $6$ 个交点。于是得到下面三条等价的刻画。

\begin{proposition}[焦点的几何刻画]
设 $g$ 是圆锥曲线，$\theta=(\sym{z}{\bar z})$ 是垂直结构。以下三条性质中的每一条都唯一确定了
$g$ 的六个焦点：
\begin{enumerate}
\item $z,\bar z$ 关于圆锥曲线的切线相交于焦点；
\item 过焦点作圆锥曲线的两条切线，它们分别经过 $z,\bar z$；
\item 圆锥曲线过 $z,\bar z$ 的四条切线交于 $6$ 个点；在这 $6$ 个点中取一对不落在同一条切线上
的点 $p,q$，令 $\theta'=(\sym{p}{q})$，则这 $6$ 个点正是 $g$ 关于垂直结构 $\theta'$ 的六个焦点。
\end{enumerate}
\end{proposition}

\begin{remark}[一般的生成点]
上面的构造并不依赖 $z,\bar z$ 的具体位置。任意给定两个点 $p,q$，都可以作出对称积
$\theta'=(\sym{p}{q})$，并写出方程
\[
\det\bigl(g^{ij}-\lambda\theta'^{ij}\bigr)=0,
\]
按完全相同的方式得到六个点。这时称这六个点为 $p,q$ 生成的焦点。第 $3$ 条性质说明，从圆锥曲线
的四条切线及其交点出发，任取合适的一对点作为新的生成点，所得的六个点保持不变。
\end{remark}

\section{准线}

前面的构造是对 $(2,0)$ 型张量进行的。同样可以对 $(0,2)$ 型张量求解
\[
\det\bigl(\theta_{ij}-\lambda g_{ij}\bigr)=0,
\qquad
\theta_{ij}=g_{ik}g_{jl}\theta^{kl}.
\]
这个方程与 \eqref{eq:focus-eq} 有相同的特征值，由它构造出的几何量与前者基本相同，只相差一次
指标的升降。

\begin{definition}[准线]
焦点的极线称为圆锥曲线的\textbf{准线}。
\end{definition}

六个焦点给出六条准线，按配对也分成三对，与三对焦点相对应。与焦点一样，通常只有一对准线落在
实仿射平面内可见的位置。

以椭圆为例。取焦点 $f_{X1}=(c,0,1)$，其极线为
\[
l_i=g_{ij}f_{X1}^j=\Bigl(\frac{c}{a^2},\ 0,\ -1\Bigr),
\]
即直线
\[
\frac{c}{a^2}x-z=0
\quad\Longleftrightarrow\quad
x=\frac{a^2}{c}.
\]
这正是解析几何中的准线 $x=a^2/c$。焦点 $f_{X2}=(-c,0,1)$ 的极线是 $x=-a^2/c$。虚焦点
$f_{Y1}=(0,\ii c,1)$ 的极线为 $\dfrac{\ii c}{b^2}y-z=0$，即
$y=-\ii b^2/c$，不在实平面内；$z,\bar z$ 的极线经过中心，也是虚的或无穷远的直线。

\section{例题：焦点、切线与准线}

\begin{example}\label{ex:focus-directrix}
设 $g$ 是椭圆，$F$ 是它的一个焦点。过 $F$ 的直线与 $g$ 交于 $A,B$ 两点，$A,B$ 处的两条切线
交于点 $H$。则 $H$ 落在 $F$ 对应的准线上，并且 $AF\perp FH$。
\end{example}

证明分两部分。

\begin{proof}[第一部分：$H$ 在准线上]
$A,B$ 是过 $F$ 的直线与 $g$ 的交点，$A,B$ 处的切线交于 $H$。由极点与极线的性质，$H$ 的极线
就是两个切点 $A,B$ 的连线，即直线 $AB$。由于 $F$ 在直线 $AB$ 上，按配极的对称性（La Hire
定理），$F$ 的极线经过 $H$。而 $F$ 的极线按定义就是 $F$ 对应的准线，因此 $H$ 在准线上。
\end{proof}

\begin{proof}[第二部分：$FA$ 与 $FH$ 共轭当且仅当它们垂直]
先说明 $FA$ 与 $FH$ 关于 $g$ 共轭。直线 $AB$ 的极点是 $H$；又 $H$ 在直线 $FH$ 上，把 $AB$
的极线 $H$ 与直线 $FH$ 相接可知 $AB$ 与 $FH$ 关于 $g$ 共轭。而 $F,A,B$ 共线，$FA$ 就是直线
$AB$，于是 $FA$ 与 $FH$ 关于 $g$ 共轭。

于是问题化为：过焦点 $F$ 的两条直线关于 $g$ 共轭，当且仅当它们垂直。这由下面两条引理给出。
\end{proof}

\begin{lemma}\label{lem:conj-harmonic}
设 $P$ 是不在圆锥曲线 $g$ 上的点，$t_1,t_2$ 是从 $P$ 引出的两条切线。则过 $P$ 的两条直线
$m,n$ 关于 $g$ 共轭（即 $g^{ij}m_in_j=0$）当且仅当 $m,n,t_1,t_2$ 构成调和线束。
\end{lemma}

\begin{proof}
共轭关系、切线与调和线束都在射影变换下保持，因此可以取 $g$ 为圆、$P$ 在一条直径上。取
\[
g:\ x^2+y^2=1,\qquad P=(d,0),\qquad d\ne0,\ \pm1 .
\]
过 $P$ 斜率为 $s$ 的直线为 $m_i=(s,-1,-sd)$，即 $y=s(x-d)$。它的极点为
\[
D^i=g^{ij}m_j=(s,-1,sd),\qquad g^{ij}=\diag\{1,1,-1\}.
\]
过 $P$ 且经过 $D$ 的直线就是 $m$ 的共轭直线，直接计算
\[
P\times D=(d,0,1)\times(s,-1,sd)=\bigl(1,\ s(1-d^2),\ -d\bigr),
\]
其斜率（表示为 $y=t(x-d)$）满足
\[
t=\frac{1}{s(d^2-1)}.
\]
这条直线与 $g$ 相切当且仅当它到原点的距离为 $1$，即
$|sd|/\sqrt{s^2+1}=1$，解得
\[
s=\pm\frac{1}{\sqrt{d^2-1}}=\pm k,\qquad k=\frac{1}{\sqrt{d^2-1}} .
\]
两条切线的斜率因此是 $\pm k$。斜率 $s$ 与其共轭直线的斜率 $k^2/s$ 的交比
\[
\Bigl(s,\ \frac{k^2}{s};\ k,\ -k\Bigr)
=\frac{(s-k)(k^2/s+k)}{(s+k)(k^2/s-k)}
=\frac{(s-k)\cdot k(k+s)/s}{(s+k)\cdot k(k-s)/s}
=-1 .
\]
故 $m,n$ 与两条切线成调和线束。
\end{proof}

\begin{lemma}\label{lem:perp-harmonic}
设 $\theta=(\sym{z}{\bar z})$ 是垂直结构，$P$ 是任一点。过 $P$ 的两条直线 $m,n$ 垂直，当且仅当
$m,n$ 与直线 $Pz,P\bar z$ 构成调和线束。
\end{lemma}

\begin{proof}
取 \eqref{eq:theta-standard} 的标准坐标系。直线 $Pz$ 与 $P\bar z$ 分别是以 $z,\bar z$ 为方向的
直线，由于 $z=(1,\ii,0)$、$\bar z=(1,-\ii,0)$，它们的斜率是 $\ii$ 与 $-\ii$。过 $P$ 的两条
直线写成 $m_i=(s,-1,b_1)$、$n_i=(t,-1,b_2)$，则
\[
\theta^{ij}m_in_j=st+1,
\]
故 $m\perp n$ 当且仅当 $st=-1$。另一方面，
\[
(s,t;\ii,-\ii)
=\frac{(s-\ii)(t+\ii)}{(s+\ii)(t-\ii)}
=-1
\]
等价于 $(s-\ii)(t+\ii)=-(s+\ii)(t-\ii)$，展开即 $st+1=-st-1$，也就是 $st=-1$。两种条件
一致。
\end{proof}

\begin{proof}[例题的完成]
设 $t_1,t_2$ 是从焦点 $F$ 引出的 $g$ 的两条切线。由焦点的构造，$F$ 是四条切线中某两条的交点，
其中一条经过 $z$、一条经过 $\bar z$；而经过 $F$ 且经过 $z$ 的切线只能是 $Fz$，经过 $F$ 且经过
$\bar z$ 的切线只能是 $F\bar z$。因此
\[
\{t_1,t_2\}=\{Fz,\ F\bar z\}.
\]
由引理 \ref{lem:conj-harmonic}，$FA$ 与 $FH$ 关于 $g$ 共轭，当且仅当 $FA,FH$ 与 $Fz,F\bar z$
构成调和线束；由引理 \ref{lem:perp-harmonic}，$FA\perp FH$ 也当且仅当 $FA,FH$ 与
$Fz,F\bar z$ 构成调和线束。两个条件等价，所以 $FA$ 与 $FH$ 共轭当且仅当 $FA\perp FH$。

前面已经证明 $FA$ 与 $FH$ 确实共轭，故 $FA\perp FH$，即 $AF\perp FH$。结合第一部分，$H$ 在
准线上。命题得证。
\end{proof}

\section{离心率的交比表示}

离心率有射影的解释。先引入一个由 $g$ 决定的交比型不变量。对两个点 $p,q$，记
\begin{equation}
\{p,q\}_g
=\frac{\bigl(g_{ij}p^iq^j\bigr)^2}
{\bigl(g_{ij}p^ip^j\bigr)\bigl(g_{ij}q^iq^j\bigr)},
\label{eq:invariant-point}
\end{equation}
只要分母不为零，它就是 $p,q$ 的规范无关的函数。对两条直线 $m,n$ 类似地记
\[
\{m,n\}_g
=\frac{\bigl(g^{ij}m_in_j\bigr)^2}
{\bigl(g^{ij}m_im_j\bigr)\bigl(g^{ij}n_in_j\bigr)} .
\]
若 $m=p^*=g_{ij}p^j$、$n=q^*=g_{ij}q^j$ 是 $p,q$ 的极线，则由 $g^{ij}m_in_j=g_{ij}p^iq^j$
等等式得 $\{p,q\}_g=\{m,n\}_g$。

\begin{lemma}\label{lem:invariant-J}
设直线 $m,n$ 分别与圆锥曲线 $g$ 交于 $a,b$ 与 $c,d$ 四个互不相同的点，$J_{ab,cd}$ 是这四点在
$g$ 上的交比，则
\begin{equation}
\{m,n\}_g=\left(\frac{J_{ab,cd}+1}{J_{ab,cd}-1}\right)^2 .
\label{eq:invariant-J}
\end{equation}
特别地，若 $p,q$ 是两点，$a,b$ 与 $c,d$ 分别是从 $p,q$ 引出的切线与 $g$ 的切点，则
$\{p,q\}_g$ 也由 \eqref{eq:invariant-J} 给出。
\end{lemma}

\begin{proof}
不妨设 $m=a\times b$、$n=c\times d$。由叉乘的恒等式
\[
(a\times b)\cdot(c\times d)=(a\cdot c)(b\cdot d)-(a\cdot d)(b\cdot c).
\]
又 $a,b$ 在 $g$ 上，$a^2=b^2=0$，故
\[
(a\times b)^2=a^2b^2-(a\cdot b)^2=-(a\cdot b)^2,
\]
同理 $(c\times d)^2=-(c\cdot d)^2$。于是
\[
\{m,n\}_g
=\frac{\bigl[(a\cdot c)(b\cdot d)-(a\cdot d)(b\cdot c)\bigr]^2}{(a\cdot b)^2(c\cdot d)^2}.
\]
由二次点列交比的公式与交比在交换下标时的变化规律，有
\[
(a\cdot c)(b\cdot d)=J^2\,(a\cdot d)(b\cdot c),\qquad
(a\cdot b)(c\cdot d)=(1-J)^2(a\cdot d)(b\cdot c),
\]
其中 $J=J_{ab,cd}$（第一式是交比平方公式，第二式来自 $J_{db,ca}=1-J_{ab,cd}$ 的平方）。代入得
\[
\{m,n\}_g
=\frac{\bigl[(J^2-1)(a\cdot d)(b\cdot c)\bigr]^2}
{\bigl[(1-J)^2(a\cdot d)(b\cdot c)\bigr]^2}
=\frac{(J+1)^2(J-1)^2}{(J-1)^4}
=\left(\frac{J+1}{J-1}\right)^2 .
\]
把上式用于极线 $m=p^*$、$n=q^*$，即得关于 $p,q$ 的结论。
\end{proof}

\begin{remark}
\eqref{eq:invariant-J} 的右端在 $J\mapsto1/J$ 下不变，因此给定 $\{p,q\}_g$ 只能把 $J$ 确定到
$J$ 与 $1/J$ 两种取法。具体取哪一个由几何连续性（例如圆时应取 $J_\theta=1$）确定。
\end{remark}

现在对椭圆 $g^{ij}=\diag\{a^2,b^2,-1\}$、$g_{ij}=\diag\{1/a^2,1/b^2,-1\}$ 逐一计算。由
\eqref{eq:zz}，
\[
z\cdot z=\frac{1}{a^2}-\frac{1}{b^2},\qquad
\bar z\cdot\bar z=\frac{1}{a^2}-\frac{1}{b^2},\qquad
z\cdot\bar z=\frac{1}{a^2}+\frac{1}{b^2},
\]
其中 $z\cdot\bar z=g_{ij}z^i\bar z^j$。因此
\begin{equation}
\{z,\bar z\}_g
=\left(\frac{1/a^2+1/b^2}{1/a^2-1/b^2}\right)^2
=\left(\frac{a^2+b^2}{c^2}\right)^2 .
\label{eq:inv-z}
\end{equation}
由 \eqref{eq:ellipse-X}，$f_{X1}=(-c,0,1)$、$f_{X2}=(c,0,1)$，于是
\[
f_{X1}\cdot f_{X2}=-\frac{c^2}{a^2}-1=-\frac{a^2+c^2}{a^2},\qquad
f_{X1}\cdot f_{X1}=f_{X2}\cdot f_{X2}=\frac{c^2}{a^2}-1=-\frac{b^2}{a^2},
\]
所以
\begin{equation}
\{f_{X1},f_{X2}\}_g
=\left(\frac{a^2+c^2}{b^2}\right)^2 .
\label{eq:inv-X}
\end{equation}
由 \eqref{eq:ellipse-Y}，$f_{Y1}=(0,\ii c,1)$、$f_{Y2}=(0,-\ii c,1)$，于是
\[
f_{Y1}\cdot f_{Y2}=\frac{c^2}{b^2}-1=\frac{c^2-b^2}{b^2},\qquad
f_{Y1}\cdot f_{Y1}=f_{Y2}\cdot f_{Y2}=-\frac{c^2}{b^2}-1=-\frac{b^2+c^2}{b^2},
\]
所以
\begin{equation}
\{f_{Y1},f_{Y2}\}_g
=\left(\frac{c^2-b^2}{c^2+b^2}\right)^2
=\left(\frac{b^2-c^2}{b^2+c^2}\right)^2 .
\label{eq:inv-Y}
\end{equation}

把 \eqref{eq:inv-z}--\eqref{eq:inv-Y} 代入 \eqref{eq:invariant-J}，并利用
$c^2=a^2-b^2$，解得相应的交比
\begin{equation}
J_\theta=\frac{b^2}{a^2}=1-e^2=\frac{\lambda_2}{\lambda_1},\qquad
J_X=e^2,\qquad
J_Y=1-e^{-2}.
\label{eq:J}
\end{equation}
例如第一个：\eqref{eq:inv-z} 给出 $\{z,\bar z\}_g=((a^2+b^2)/c^2)^2$，而
$(J+1)/(J-1)$ 在 $J=b^2/a^2$ 时等于 $(a^2+b^2)/(b^2-a^2)=-(a^2+b^2)/c^2$，平方后与
\eqref{eq:inv-z} 一致，故 $J_\theta=b^2/a^2=1-e^2$。另外两式同法可得，其中 $J_X=e^2$ 由
$(e^2+1)/(e^2-1)=-(a^2+c^2)/b^2$ 得到，$J_Y=1-e^{-2}$ 由
$(J_Y+1)/(J_Y-1)=(b^2-c^2)/(b^2+c^2)$ 得到。这里 $J_\theta$ 是过 $z,\bar z$ 的四条切线的
四个切点所对应的交比，$J_X$ 与 $J_Y$ 是把 $z,\bar z$ 分别换成 $f_{X1},f_{X2}$ 与
$f_{Y1},f_{Y2}$ 后所得的同类交比。于是离心率完全由这三个交比决定。

\subsection{特殊情形与连续变化}

由 \eqref{eq:J} 可以读出几个特殊位置。
\begin{itemize}
\item $J_\theta=-1$ 时四线成调和线束，此时 $1-e^2=-1$，$e=\sqrt2$，圆锥曲线是等轴双曲线。
\item $J_\theta=0$ 时成立 $b^2=0$，$e=1$。此时 $\{z,\bar z\}_g=1$ 所对应的四条切线中有两条
重合，有一个焦点落在圆锥曲线上，圆锥曲线与无穷远直线相切，即抛物线。
\item $e=0$ 时 $\lambda_1=\lambda_2$，圆锥曲线是圆。此时 $z,\bar z$ 都落在圆锥曲线上，从它们
引出的切线各退化为一条；四个有限焦点同时重合到圆锥曲线的中心，也就是无穷远直线的极点。
\end{itemize}
相应地，$J_X=e^2$ 在圆时为 $0$，在抛物线时为 $1$，在等轴双曲线时为 $2$；$J_Y=1-e^{-2}$ 在
圆时为无穷（即该对焦点的连线退化），在抛物线时为 $0$，在等轴双曲线时为 $1/2$。

椭圆的离心率变化时，六个焦点的位置可以连续地描述。设 $a>b>0$，四个有限焦点是
$(\pm c,0)$ 与 $(0,\pm\ii c)$，其中 $c^2=a^2-b^2$。当 $a$ 与 $b$ 逐渐接近时，$c$ 减小，
$x$ 轴上的两个实焦点与 $y$ 轴上的两个虚焦点同时向中心靠拢；当 $a=b$ 时 $c=0$，四个焦点重合
于中心，此时圆锥曲线是圆，而 $z,\bar z$ 落在圆锥曲线上。若继续压缩，使 $b>a$，则原来的两个
实焦点变为虚焦点、原来的两个虚焦点变为实焦点，四个焦点重新分开：此时长轴转到 $y$ 轴，实焦点
位于 $y$ 轴。整个过程中两个无穷远焦点 $z,\bar z$ 的位置不变。

\section{注：共焦点圆锥曲线族}

最后补充关于共焦点圆锥曲线族的一个观察。

\begin{remark}[共焦点族]
设 $g,h$ 是两条圆锥曲线。方程
\[
\det(g-\lambda h)=0
\]
的根给出与 $g,h$ 的焦点有关的射影数据；当 $g,h$ 共焦点时，这个方程出现重根。

另一方面，焦点由四条切线决定。设 $g$ 及其过 $z,\bar z$ 的四条切线给定，则所有与这四条切线都
相切的圆锥曲线 $h$ 都与 $g$ 有相同的焦点，它们构成一个单参数的圆锥曲线族，即共焦点圆锥曲线
族。因此共焦点族既可以用方程 $\det(g-\lambda h)=0$ 从代数上刻画，也可以用四条公切线的相切条件
从几何上刻画。这两种刻画互相等价，可作为处理共焦点圆锥曲线的一种方法。
\end{remark}
